\section{Canonical normal forms and the Euler principle}
\label{sec:normal}

Let
\[
 \Rcal=\Q[u_0,u_1,u_2,\ldots],\qquad
 \D u_j=u_{j+1},\qquad \wt(u_j)=j+1.
\]
The factor degree of a monomial is its ordinary total degree.
A monomial is \emph{functional} if it is a power of $u_0$, including $1$,
or if its largest positive derivative index occurs at least twice.
Denote the rational span of these monomials by $\mathscr F$.

\begin{proposition}[Canonical decomposition]\label{prop:normal}
Every $P\in\Rcal$ has a unique decomposition
\begin{equation}\label{eq:decomp}
 P=\NF(P)+\D\Prim(P),\qquad
 \NF(P)\in\mathscr F,\qquad [1]\Prim(P)=0.
\end{equation}
Both maps are rational and linear. They preserve factor degree; on a
homogeneous input of weight $m$, $\NF(P)$ has weight $m$ and $\Prim(P)$
has weight $m-1$.
\end{proposition}
\begin{proof}
Suppose that a monomial $M$ has a unique largest derivative $u_r$,
$r\ge1$. Replace that factor by $u_{r-1}$, obtaining $T$, and let $c$
be the multiplicity of $u_{r-1}$ in $T$. Then
\[
 M=\frac1c\D T+\left(M-\frac1c\D T\right).
\]
The remaining monomials have strictly smaller largest derivative index.
Repeated application terminates and constructs both maps. The asserted
gradings are preserved at every step.

For uniqueness, let a nonconstant $H$ have largest index $r$. In $\D H$
the term $u_{r+1}\partial_{u_r}H$ is nonzero and consists entirely of
monomials with unique largest index $r+1$. It cannot be canceled by the
other terms of $\D H$, which have largest index at most $r$, or by a
functional polynomial. Thus $\D H\notin\mathscr F$. This proves
uniqueness once the constant term of $H$ is fixed.
\end{proof}

The decomposition is a direct sum of vector spaces, not of algebras.
In particular, $\NF$ is not multiplicative: $\NF(u_1)=0$ while
$\NF(u_1^2)=u_1^2$. The constructive proof gives a complete endpoint
primitive, and the supplied rational program implements that proof.

\subsection{The inverse-Gamma Bell hierarchy}
Define
\begin{align}
 A_0&=1,& A_{m+1}&=u_0A_m-\D A_m,\label{eq:Bell-recur}\\
 \Acal(t)&=\sum_{m\ge0}A_m\frac{t^m}{m!}
 =\exp\left(\sum_{r\ge0}\frac{(-1)^ru_rt^{r+1}}{(r+1)!}\right),
 \label{eq:Bell}\end{align}
and put $N_m=\NF(A_m)$ and $H_m=\Prim(A_m)$.
Formula~\eqref{eq:Bell} follows either by differentiating the exponential
or by the formal Taylor formula. Under $u_r=\psi^{(r)}(x)$,
\begin{equation}\label{eq:Gamma-realization}
 A_m(x)=(-1)^m\Gamma(x)\frac{\mathrm d^m}{\mathrm dx^m}\frac1{\Gamma(x)},
 \qquad \Acal(t,x)=\frac{\Gamma(x)}{\Gamma(x-t)}.
\end{equation}
These identifications use $\psi=\Gamma'/\Gamma$ and its derivatives
\cite{DLMF5}. For every fixed $m$ they are ordinary analytic identities
away from Gamma poles.

Translation $u_0\mapsto u_0+c$ commutes with $\D$ and preserves
functional monomials. By uniqueness it commutes with $\NF$. Consequently
$\NF\Acal(t)$ has an Appell factor $e^{u_0t}$. A second useful identity is
\begin{equation}\label{eq:centered}
 \NF\Acal(t)=\NF\exp\left(
 \sum_{j\ge0}\frac{u_{2j}t^{2j+1}}{2^{2j}(2j+1)!}\right).
\end{equation}
Indeed the exponential before normalization on the right is
$e^{t\D/2}\Acal(t)$. Every coefficient of
$(e^{t\D/2}-1)\Acal(t)$ is a total derivative. The exponent in
\eqref{eq:centered} has odd weights only, so every retained monomial has
factor degree congruent to its weight modulo two. For
$u_0^a u_1^k u_2^b u_3^n$ this says that $k+n$ must be even.

\subsection{Euler derivatives and finite jets}
Define
\[
 \Euler(P)=\sum_{r\ge0}(-\D)^r\partial_{u_r}P.
\]
Every sum is finite on a polynomial. The commutator relation
$\partial_{u_r}\D=\D\partial_{u_r}+\partial_{u_{r-1}}$, with the
last term zero at $r=0$, gives $\Euler\D=0$ by cancellation.

Let $\sigma_d$ set all $u_r$ with $r>d$ to zero and define
$\delta_d=\sum_{r=0}^{d-1}u_{r+1}\partial_{u_r}$.

\begin{lemma}[Finite-jet variational identity]\label{lem:finitejet}
For every $P\in\Rcal$,
\begin{equation}\label{eq:finitejet}
 \sigma_d\Euler(P)=
 \sum_{r=0}^d(-\delta_d)^r\partial_{u_r}\bigl(\sigma_d\NF(P)\bigr).
\end{equation}
Furthermore,
\begin{equation}\label{eq:Euler-A}
 \Euler\Acal(t)=\int_0^t e^{s\D}\Acal(t)\dd s.
\end{equation}
The integral in \eqref{eq:Euler-A} is coefficientwise formal.
\end{lemma}
\begin{proof}
Since $\Euler\D=0$, replace $P$ by $\NF(P)$ on the left of
\eqref{eq:finitejet}. A normal monomial removed by $\sigma_d$ has at
least two factors of its highest index, which exceeds $d$. One partial
derivative removes at most one such factor; subsequent applications of
$\D$ preserve a factor of index greater than $d$. Hence its specialization
vanishes. For retained monomials, the ideal generated by
$u_{d+1},u_{d+2},\ldots$ is $\D$-stable, and specialization changes
$\D$ to $\delta_d$. This proves \eqref{eq:finitejet}.

By \eqref{eq:Bell},
$\partial_{u_r}\Acal(t)=(-1)^rt^{r+1}\Acal(t)/(r+1)!$. Therefore
\[
 \Euler\Acal(t)=\sum_{r\ge0}\frac{t^{r+1}}{(r+1)!}\D^r\Acal(t),
\]
which is the formal integral in \eqref{eq:Euler-A}.
\end{proof}

Normal reduction must precede truncation. For example,
$\sigma_1(u_0u_2)=0$, whereas
$\sigma_1\NF(u_0u_2)=-u_1^2$.
The finite-jet lemma, not commutation of these operations, is the reason
one can obtain a closed low-derivative equation.
